\section{Keep one working source through the observations}
Copy \path{drafts/cbor-construction/main.asm} to a working file. Keep that working
file as each region is completed. The caller has modes for head reading, generic
skipping, complete decoding, exact sizing, encoding and the connected pipeline. Two
additional modes call the illustrated private head-size and head-writing helpers.
Their labels and register choices support guided construction; an equivalent complete
solution need only meet the five public entry contracts.

From the construction directory, run
\texttt{python3 prepare.py s01-1 --source working.asm --output observation.asm}.
Use the exact recipe named in the step and a fresh output filename. Preparation
changes fixture bytes and scalar controls, while preserving the entire instruction
section, including comments and instructions already written. Continue editing the
working file; the prepared file supplies the selected observation data.

Every region starts with a safe branch to a development return. During construction,
status four with a zero count means that an unfinished region was reached. It is
outside the completed statuses zero through two. A later unfinished body may follow
valid earlier work, including partial encoding or publication after full input
validation. Keep development status distinct from a completed format or capacity
failure. When calling another unfinished routine, propagate its status before using
its missing result.

Stop before the named checkpoint instruction. Each supplied after-label owns a
\texttt{nop}, so adjacent checkpoints have distinct addresses. A branch taken to a
handler may bypass an after-label; the recipe names the observation on its intended
path. A direct \texttt{function\_return} checkpoint corresponds to the selected
caller's return label. \texttt{results\_ready} is before the caller stores its result
slots and before the exit syscall replaces the live status register.

\section{Use the supplied frames as private working storage}
The walker has an eighty-byte frame, with its return address at zero and saved
\texttt{s0} through \texttt{s6} at offsets eight through fifty-six. The decoder uses
288 bytes, saving its return address and \texttt{s0} through \texttt{s7} at offsets
zero through sixty-four. The original input count is retained at eighty and the nested
head helper's return at eighty-eight. Its candidate starts at ninety-six: counter at
ninety-six, lengths at 104 and 108, identifier at 112 and sample slots at 144. These
locations are inside the private frame, not the caller's record.

The bound frame is sixty-four bytes, the private writer frame thirty-two and the
encoder frame eighty. Their supplied saves and restores retain the registers used
by this illustrated organization. The head reader and size helper are leaves and
need no frame for their temporary-register work. Do not restore a frame on an early
count failure that occurred before its allocation.

Caller data controls are full-width unsigned bit patterns. A decimal value above the
signed maximum is emitted with the same bits as a negative assembly literal; it is
still an unsigned count or argument under its declared contract. Bound and encode
recipes seed a flat record. Decode and pipeline recipes start that record as sentinel
bytes, so an unchanged output cannot resemble a prepared successful result. Step 55
uses two separate recipes with different caller modes to compare two failure contracts.
